% -------------------------------------------------------
\section{Introduction}
\label{sec:introduction}

The deployment of autonomous multi-agent systems in high-stakes physical environments — irrigation networks, autonomous vehicles, distributed manufacturing, healthcare logistics — has exposed a structural gap in the prevailing approach to AI accountability. Modern multi-agent architectures delegate consequential decisions to autonomous actors whose authority to act is assumed rather than enforced. When something goes wrong, the accountability chain is recovered through post-hoc investigation, not architectural guarantee. We call this the \textit{black box accountability problem}: autonomous actions that no human reviewed before execution, cannot be explained after the fact with fidelity, and for which no human bears unambiguous legal or operational responsibility.

The problem is not incidental. It is architectural. The dominant multi-agent frameworks treat human oversight as an optional layer — a governance document, an audit log, a review process that can be added or removed without affecting the system's core functionality. This treatment reflects an organizational intuition about accountability rather than an engineering constraint. In this paper, we argue that accountability without architectural enforcement is not accountability at all: it is documentation that happens to describe actions whose real-world consequences have already occurred.

% -------------------------------------------------------
\subsection*{Why Autonomous Multi-Agent Systems Require a Mandatory Bail-Out Mechanism}
\label{sec:intro:why-mandatory}

Consider the structure of a typical autonomous multi-agent deployment. A root orchestrator spawns sub-agents. Sub-agents may spawn further sub-agents. Each agent performs inference, invokes tools, makes decisions, and drives physical or digital actuators. The spawning hierarchy creates a delegation chain: a human at the top authorizes objectives; intermediate agents translate those objectives into operational decisions; leaf agents execute. The delegation chain is the accountability chain — in principle.

In practice, the chain is broken at the first unresolvable exception. When a leaf agent encounters a condition it cannot handle — an out-of-range sensor reading, a safety-envelope violation that its own simulation cannot resolve, a decision whose implications exceed its legal or operational authority — it does not escalate to a human. It either fails silently, applies a local heuristic that may or may not be appropriate, or propagates the error upward in a form that no human can act on coherently. The exception is absorbed by the machine hierarchy. The human at the top of the chain is never informed. The accountability chain terminates not at a human, but at an unresolved state that has already produced physical consequences.

This failure mode is not hypothetical. In deployed autonomous systems, it manifests as: an autonomous vehicle that continues into a condition its safety model cannot evaluate rather than requesting driver intervention; an industrial robotic arm that applies excessive force rather than entering a safe hold state because the force sensor anomaly was filtered out by a confidence threshold; an irrigation controller that continues a watering cycle through a pipe fracture rather than executing a bailout because the fracture signature was ambiguous in the sensor fusion model. In each case, the architecture provided no path for an unresolvable condition to reach a human decision-maker before the consequential action was taken.

The mandatory bail-out mechanism we propose in this paper is not a ``big red button.'' It is an architectural protocol embedded in every node of the spawning hierarchy, such that any unresolvable condition at any depth propagates asynchronously to a Human Accountability Anchor, bypassing all intermediate machine actors, with complete diagnostic context. The protocol is structurally mandatory: it cannot be suppressed by a compromised node, cannot be bypassed by a machine actor, and cannot be silenced by a network partition without detection. It is the direct conversion of Dijkstra's separation-of-concerns principle into an enforcement mechanism for human accountability.

The structural mandatory property is the essential contribution. Existing approaches to exception handling in multi-agent systems treat escalation as a design choice to be made at deployment: if the system designer anticipated a condition, they may have included a human review step; if they did not, the condition is handled by the machine hierarchy alone. The Bailout Protocol removes this dependence on anticipation. Any condition that the machine layers cannot resolve with certainty — whether anticipated or not — triggers the same mandatory escalation. The protocol does not enumerate failure modes; it defines a single resolution path for all unresolvable states.

% -------------------------------------------------------
\subsection*{What Happens When Systems Do Not Have One}
\label{sec:intro:what-happens-without}

The absence of a mandatory bail-out mechanism does not produce neutral outcomes. It produces a specific failure pattern we call \textit{autonomous drift}: the progressive extension of machine authority beyond its intended bounds, without human review, until a failure occurs that forces human intervention under suboptimal conditions.

Autonomous drift is enabled by the confidence-threshold model that dominates existing HITL frameworks. In this model, a human reviews an agent's action only when the agent's confidence falls below a configured threshold. The model assumes that an agent's self-reported confidence is a reliable signal of whether the action is correct with respect to the user's actual objectives. This assumption fails systematically for large language model--based agents, whose confidence scores are calibrated against the model's training objective (next-token prediction), not against the operational objective (correct action in context). A model can express high confidence while being wrong about what the user intended — and the confidence-threshold model will permit the action to proceed without human review.

The consequences of autonomous drift are compounding. When an agent acts without human review and the action produces an acceptable outcome, the system reinforces the pattern: this agent handled this condition autonomously, no escalation was necessary, the confidence threshold was appropriate. When the same agent encounters a similar condition and acts similarly, the pattern strengthens. Over successive cycles, machine authority extends into territory that was never authorized, never reviewed, and never audited. The accountability chain has not been broken — it has been quietly replaced by a machine chain that no human is monitoring.

The eventual failure is then catastrophic in a specific sense: it occurs under conditions the human never reviewed, about whose context they have no diagnostic record, and whose resolution requires intervention in a physical state that has already diverged significantly from the intended state. The human is not a backup safety layer. They are a last-resort responder called in after damage has occurred. This is the inverse of fail-safe design. The system has failed dangerously, and the human is responding to the failure rather than preventing it.

The regulatory environment is beginning to respond to this pattern. The EU AI Act and the NIST AI Risk Management Framework both mandate human oversight for high-risk AI systems as a compliance obligation, not a best practice. These mandates presuppose an architecture capable of satisfying them. A system that relies on confidence thresholds and post-hoc audit logs cannot satisfy the mandate because it cannot guarantee that every unresolvable state reached a human before consequential action was taken. The Bailout Protocol is designed to be the architectural substrate that makes these regulatory mandates materially satisfiable.

% -------------------------------------------------------
\subsection*{The Core Insight: Accountability Must Have an Architectural Fallback}
\label{sec:intro:core-insight}

The core insight of this paper is that accountability cannot be a property of a system's governance documents — it must be a property of its architecture. Specifically: human accountability must be implemented as a structurally mandatory fallback, not as a configurable threshold or an organizational policy.

This distinction matters because governance documents and architectural constraints have fundamentally different failure modes. A governance document can be updated, overridden, or ignored by a system operator. An architectural constraint cannot be bypassed without modifying the system's structure, and any modification is logged. When accountability is implemented as a governance document, it exists at the discretion of the system operator. When accountability is implemented as an architectural fallback, it exists independent of any individual operator's choices — including the choices of the humans nominally responsible for the system.

The \bxthree Framework implements this insight through a three-layer separation of concerns: the \purpose layer, which owns all human accountability anchors; the \bounds layer, which proposes actions within authorized constraints; and the \fact layer, which enforces deterministic physical and operational constraints at the actuator interface. The Bailout Protocol is the protocol that connects these layers: it guarantees that any state the \bounds layer cannot resolve within its current authority escalates to a \purpose layer anchor, bypassing all intermediate \bounds and \fact actors, with complete diagnostic context drawn from the Forensic Ledger.

The architectural fallback property has a precise meaning in this context. A fallback is not a ``backup plan'' that operates when the primary plan fails. It is a mandatory termination condition that the system is structurally required to reach before certain classes of action can be taken. In the Bailout Protocol, the termination condition is a Human Accountability Anchor — a named human with organizational authority — and the mandatory property is enforced by the immutability of parent pointers in the spawning hierarchy. No agent can modify its own parent pointer, and no machine actor can redirect an escalation away from its designated Human Accountability Anchor without this redirection being detected by the Forensic Ledger.

This is the architectural guarantee that no existing multi-agent framework provides. Existing systems implement exception handling as a design choice, confidence thresholds as a configurable parameter, and human oversight as an organizational policy. The \bxthree Framework implements human oversight as a mathematical consequence of immutable parent pointers, depth-bounded spawning, and the mandatory termination property of the Bailout Protocol. The accountability chain does not rely on the good judgment of system operators. It is structurally enforced.

% -------------------------------------------------------
\subsection*{Paper Roadmap}
\label{sec:intro:roadmap}

The remainder of this paper is organized as follows. Section~\ref{sec:background} reviews the foundational threads converging in the Bailout Protocol: the principle of separation of concerns as an accountability architecture, the problem of accountability in multi-agent AI systems, human-in-the-loop design and its structural limitations, and fail-safe design in safety-critical systems. Section~\ref{sec{design}} presents the formal design of the Bailout Protocol, including its triggering conditions, escalation path, trigger lifecycle state machine, and the conditions under which containment is achieved versus escalated. Section~\ref{sec:correctness}} establishes the protocol's correctness properties: termination in bounded depth, integrity of the escalation path against tampering, and drift detection through forensic ledger comparison. Section~\ref{sec:forensic-ledger}} describes the Forensic Ledger as the audit substrate on which the Bailout Protocol operates, including the three-plane recording model and the hash-chained append-only structure. Section~\ref{sec:limitations}} acknowledges the protocol's limitations — human response latency, adversarial conditions, scalability at deployment depth, and exploitation attack vectors — and proposes defined directions for future work. Section~\ref{sec:conclusion}} concludes with a discussion of the protocol's relationship to the \bxthree Framework and the broader implications for mandatory human accountability in autonomous systems.

\end{document}
